% !PRINTED 194
\weberpage{194}
The behaviour of $H$, $D$, $Q$ under linear transformation follows from \S.\,60, (2), (3); if $H'$, $D'$, $Q'$ are the corresponding constructions for a transformed form, one has
\weq{8}{H' = r^2\,H,\quad D' = r^6\,D,\quad Q' = r^3\,Q.}

We can use the covariants to transform the cubic form to a normal form which at the same time gives the solution of the cubic equation.

We choose the normal form
\weq{9}{f(x,y) = F(\xi,\eta) = \xi^3 + \eta^3,}
in which $\xi$, $\eta$ are linear functions of $x$, $y$; a form which, as will appear presently, can always be produced when $D$ does not vanish, that is, when the equation $f = 0$ does not have two equal roots. The solutions of $f = 0$ are then obtained from the linear equations
\[
\xi + \eta = 0,\quad \xi + \varepsilon\,\eta = 0,\quad \xi + \varepsilon^2\,\eta = 0,
\]
in which $\varepsilon$ is a third root of unity.

To form the expressions $H'$, $D'$, $Q'$ for the normal form (9), we have to set $a_0' = a_3' = 1$, $a_1' = a_2' = 0$, and obtain
\[
H' = 9\,\xi\,\eta,\quad D' = -\,27,\quad Q' = 27\,(\xi^3 - \eta^3).
\]

If the normal form (9) is derived from the general form $f(x,y)$ by linear transformation, then by (8) we obtain
\weq{10}{
\begin{aligned}
9\,\xi\,\eta &= r^2\,H \\
-\,27 &= r^6\,D \\
27\,(\xi^3 - \eta^3) &= r^3\,Q \\
\xi^3 + \eta^3 &= f.
\end{aligned}
}
From this it follows that
\[
r^6 = \frac{-\,27}{D},\quad 2\,\xi^3 = \frac{r^3\,Q}{27} + f,\quad 2\,\eta^3 = -\,\frac{r^3\,Q}{27} + f,
\]
hence, if for $r^3$ the value $-\,9 : \sqrt{-3\,D}$ from the first equation is substituted,
\weq{11}{
\begin{aligned}
6\,\xi^3\,\sqrt{-3\,D} &= +\,Q + 3\,\sqrt{-3\,D}\,f \\
6\,\eta^3\,\sqrt{-3\,D} &= -\,Q + 3\,\sqrt{-3\,D}\,f,
\end{aligned}
}
from which $\xi$ and $\eta$ are determined by two cube roots, of which, by the first equation (10), the one is determined by the other.
